% All \usepackage commands should now be in main.tex

% Probability and Statistics
\DeclareMathOperator*{\PrExp}{\mathbb{E}} % Operator for expectation (if used this way)
\DeclareMathOperator{\Prob}{\Pr}           % Probability operator
\DeclareMathOperator{\Exp}{\mathbb{E}}     % Expectation (often used as \mathbb{E}[X])
\DeclareMathOperator{\Var}{\mathrm{Var}}   % Variance operator
\newcommand{\ass}{\operatorname{assoc}}

% Sets and Delimiters
\newcommand{\set}[1]{\left\{#1\right\}}
\newcommand{\abs}[1]{\left|#1\right|}
\newcommand{\norm}[1]{\left\|#1\right\|}
\newcommand{\floor}[1]{\lfloor #1 \rfloor}
\newcommand{\ceil}[1]{\lceil #1 \rceil}

% Common delimiters with \left and \right
\newcommand{\bra}[1]{\left(#1\right)}
\newcommand{\sbra}[1]{\left[#1\right]}
\newcommand{\cbra}[1]{\left\{#1\right\}}
\newcommand{\abra}[1]{\left\langle#1\right\rangle}

% Fields/Number Systems
\newcommand{\F}{\mathbb{F}}
\newcommand{\R}{\mathbb{R}}
\newcommand{\N}{\mathbb{N}}
\newcommand{\Z}{\mathbb{Z}}
\newcommand{\C}{\mathbb{C}}

% Common notations
\newcommand{\eps}{\varepsilon}
\newcommand{\poly}{\mathrm{poly}} % for polynomial complexity
\newcommand{\zo}{\{0,1\}}
\newcommand{\negl}{\mathrm{negl}} % for negligible functions
\newcommand{\defeq}{\stackrel{\text{def}}{=}} % Defined as

% Specific to your paper: Algorithms and Functions (using \DeclareMathOperator)
\DeclareMathOperator{\DTdepth}{DTdepth}
\DeclareMathOperator{\CDTdepth}{CDTdepth}
\DeclareMathOperator{\CDT}{CDT}
\DeclareMathOperator{\Pivot}{Pivot}
\DeclareMathOperator{\BalanceZero}{0-Balance}
\DeclareMathOperator{\BalanceOne}{1-Balance}
\DeclareMathOperator{\PathVars}{PathVars} % Function returning variables on a path
\DeclareMathOperator{\assoc}{Assoc}
\newcommand{\DT}{\mathsf{DT}}

% Specific to your paper: Symbols and Structures
\newcommand{\unassigned}{\star} % For unassigned variables in restrictions (used for rho(x_i) = \star)
\newcommand{\ot}{\leftarrow}

% Calligraphic for trees and events
\newcommand{\calE}{\mathcal{E}}
\newcommand{\calA}{\mathcal{A}}
\newcommand{\calB}{\mathcal{B}}
\newcommand{\calT}{\mathcal{T}}
\newcommand{\cR}{\mathcal{R}}

% Bold italic for sequences/vectors (to avoid conflict with random variable notation)
\newcommand{\seqvar}{\bm{x}}         % Sequence of variables (e.g., path in DT)
\newcommand{\seqassign}{\bm{a}}       % Sequence of assignment bits (e.g., path in DT)
\newcommand{\seqvarv}{\bm{x}_v}       % Sequence of variables under node v for subformula F_v
\newcommand{\pv}{\mathrm{PathVars}} 
\newcommand{\Leaves}{\operatorname{Leaves}}

\newcommand{\condF}{\pv_0(F,\rho,a) = (x_{i_1}, \dots ,x_{i_t})}
\newcommand{\condLeft}{\pv_0(F_1,\rho,a_{\le s}) = (x_{i_1}, \dots ,x_{i_s})}
\newcommand{\condRight}{\pv_0(F,\rho,a_{>s}) = x_{i_{s+1}}, \dots x_{i_t}} 
\newcommand{\condAssoc}{Assoc_1(F_1,\rho,a_{\le s}) = b} 
\newcommand{\cE}{\mathcal{E}} 
\newcommand{\varseq}[2]{x_{i_{#1}}, \dots x_{i_{#2}}}


% Placeholder for measure function (now defined as mu)
\newcommand{\Lfunc}{\mathrm{L}} % Using mu for measure. Use as \Lfunc(arg) or \Lfunc_{subscript}

% For comments/notes (xcolor package is loaded in main.tex)
\newcommand{\todo}[1]{{\color{red}\textbf{TODO: #1}}}
\newcommand{\tm}[1]{{\color{blue}[TM: #1]}}
\newcommand{\ph}[1]{{\color{purple}[PH: #1]}} % Changed to purple for distinction
\newcommand{\jr}[1]{{\color{green}[JR: #1]}}
\def\hastad{H{\aa}stad}

\def\demorgan{De Morgan}